\documentclass[10pt,a4paper]{article}

\usepackage[margin=2.4cm]{geometry}
\usepackage[T1]{fontenc}
\usepackage[utf8]{inputenc}
\usepackage{times}
\usepackage{amsmath}
\usepackage{amssymb}

\makeatletter
\renewcommand{\section}{\@startsection{section}{1}{0pt}%
  {1.3em}{0.5em}{\large\bfseries}}
\makeatother
\setlength{\parskip}{0.45em}
\setlength{\parindent}{0pt}
\pagestyle{plain}

\begin{document}

\begin{center}
{\large\bfseries Response to Reviewers}\\[0.3em]
{\itshape The Energy Cost of Privacy in Federated Learning}\\[0.15em]
cSIS 2026 --- Track 1, Green \& Intelligent Computing
\end{center}

\hrule height 0.4pt
\medskip

We thank both reviewers for their careful reading. The second reviewer's
comments led us to re-derive several numbers directly from the raw logs; the
most consequential outcome is a correction to the privacy guarantee (Point 3 of
Reviewer 2), which we now report as a finding in its own right. All tables,
figures and in-text numbers in the revision are generated by a single analysis
script from one measurement sweep per dataset, and no figure has been edited by
hand.

\section*{Reviewer \#1}

We thank the reviewer for the accurate summary of our contributions and for
recognizing the experimental controls, which were central to the study's
design.

\textbf{Page limit.} The revision fits within six pages. The space came from
removing duplication between Sections IV and V and from shortening the
captions; the carbon figure was also dropped, since Table III carries the same
six numbers. No experimental result was removed for space.

\textbf{Keywords.} The heading now reads ``Keywords''.

\textbf{Title.} The title is now set in bold.

\textbf{Captions.} All figure and table captions now state only what is shown,
the units and the averaging. Interpretive material has moved into the body
text.

\section*{Reviewer \#2}

\textbf{1. Two baselines and two noise floors.} We thank the reviewer for
identifying this. The submitted manuscript combined two measurements of the
non-private condition: the randomized sweep that all other conditions come
from, and a separate later re-run of that condition alone, used in Fig.~2 and
the opening of Section~IV. Only the sweep survives --- the re-run's files were
overwritten by a subsequent pass and cannot be regenerated --- so we can defend
only the sweep, and the revision uses it throughout. The two differ by
$3.7\%$, inside the noise floor, so nothing downstream changes qualitatively.
Every number in Tables~I--III, Figs.~1--2 and the text now comes from that
sweep: baseline $13.98$~J/round, mean overhead $+51\%$, range $+43.4\%$ to
$+62.2\%$, board-power noise floor $7.80\%$. The $14.5$~J / $+45\%$ / $8.2\%$ /
$9.46\%$ set is withdrawn.

On $\varepsilon$-dependence we also changed the instrument, because the
submitted comparison was mismatched twice over. The $12.5\%$ spread is a spread
in net energy per round, while the $7.80\%$ floor is a spread in board power;
and a floor bounds the seeds of \emph{one} condition, whereas the claim
concerns the spread of \emph{five condition means}. Each floor is now measured
on the quantity it bounds --- on energy per round it is $12.0\%$, not $7.8\%$
--- and the ordering claim is tested rather than compared to a floor: permuting
the seed labels across the five private budgets gives $F = 1.46$, $p = 0.27$ on
HAR and $F = 0.46$, $p = 0.79$ on CIFAR-10. The contrast the reviewer raises,
$+43.4\%$ at $\varepsilon = 1$ against $+62.2\%$ at $\varepsilon = 4$, is not
resolvable: the three seeds of the $\varepsilon = 4$ condition span $20.9$ to
$24.1$~J/round on their own, and the overheads are not ordered in $\varepsilon$.
Section~IV-A is retitled accordingly and claims no resolvable dependence on
$\varepsilon$, and explicitly no equality. Across budgets, per-round energy
varies by at most $1.13\times$, against $1.28\times$ to unbounded for the
round-count term.

\textbf{2. Numbers not reproducible from the tables.} The 0.50-target range,
the 0.70 ratio and the 0.85 entry in the submitted text were stale and have
been regenerated ($1.28$--$1.79\times$, $5.7\times$, $839$~J over $57.0$
rounds). The per-round discrepancies the reviewer identifies had a different
cause: the $2477$~J post-peak figure and the $29$~J / $64$~J ablation figures
were gross energy, while Table~I reports net of idle. Gross energy at
$\varepsilon = 0.5$ is $28.8$~J/round, which over 83 rounds gives
$\approx\!2390$~J; the two ablation cells are $30.5$ and $64.7$~J/round gross
against $22.0$ and $48.6$ net. The revision reports every energy figure net of
idle, labels any gross figure explicitly, and generates all in-text numbers
from the same arrays as the tables. We also note in Section~IV-E that the
ablation sweep's own baseline sits about $2$~J/round above Table~I's, so
ablation energies are compared only within that sweep.

\textbf{3. Privacy guarantee and step count.} We agree on both points.
(a)~Per-sample clipping and noising within each client yields example-level DP,
enforced locally against an untrusted server and accounted per client over the
whole run; it is not client-level DP. On HAR it protects individual windows of
a participant's data, not the participant's participation. The guarantee is
relabelled throughout (abstract, Sections~I, II-D, III-B, V-E, conclusion), and
[2] and [11] are now cited as client-level mechanisms, stronger than ours; the
per-round overhead we measure is the per-sample gradient work that both need,
so the energy result is orthogonal to the choice between them. (b)~The reviewer
is right that calibrating $\sigma$ on expected participation does not bound the
realized guarantee. We recomputed the realized $\varepsilon$ of every client
from its logged participation with the same RDP accountant. Clients joined
between 33 and 64 rounds against an expected 50, and the worst realized
$\varepsilon$ exceeds the nominal budget in every private run --- by $4$--$14\%$
on HAR ($0.570$ for nominal $0.5$, $9.11$ for nominal $8$) and by up to $22\%$
on CIFAR-10. Table~I now reports the worst realized $\varepsilon$ per
condition, the sweep is described as a sweep over nominal budgets, and the
claim that the reported $\varepsilon$ is an upper bound has been removed;
Section~III-B states the remedy, which is to calibrate against
$R\,E\,\lceil n/B \rceil$ or cap participation per client, and the audit script
reports that bound directly. The energy measurements are unaffected, since they
depend only on $\sigma$. We now present this as a second accounting hazard
alongside the calibration-timing hazard: both are silent in the resulting
numbers, and only a replay of the logs exposes either.

\textbf{4. The 83\% and 21.4$\times$ figures.} We agree, and a per-seed
analysis prompted by this comment goes further. On HAR at $\varepsilon = 0.5$
the three seeds peak at rounds 16, 54 and 65, and the mean curve shows no early
peak. The HAR stopping-round claim is therefore not supported and has been
withdrawn, together with the $83\%$ figure, which, as the reviewer notes,
restated the schedule. On CIFAR-10 the per-seed check does support the effect,
but only at the two tightest budgets: at $\varepsilon = 0.5$ the decline is
$0.026$ against a $\pm 0.010$ across-seed spread with the seeds peaking at
rounds 8, 8 and 9, and at $\varepsilon = 1$ it is $0.044$ against $\pm 0.014$
with the seeds between rounds 10 and 16. At $\varepsilon \ge 2$ the declines
fall inside their spreads. We therefore retain the stopping round on CIFAR-10
at $\varepsilon \le 1$ and withdraw the monotone ordering of the peak round,
which only two budgets resolve. Section~IV-C now leads with the peak's position
rather than the share of joules after it, and reports a peak only where the
decline clears that condition's spread and the seeds agree on where it falls.
We also introduce a separate accuracy noise floor, since the energy floor had
been applied to accuracy differences. On carbon, we agree that $21.4\times$ is
the grid-intensity ratio and would apply equally to non-private training. It
has been removed from the abstract, and Section~IV-F now presents per-grid
carbon attributable to privacy as context for federated deployment, where no
operator chooses the grid.

\textbf{5. References.} We thank the reviewer for the verification. The
truncation to one author was the bibliography style's default name limit rather
than missing data; the style now lists up to six authors before ``et al.'' as
IEEE requires. Volume, issue and page numbers have been added where they exist,
and the DOI for Aroussi et al.\ is
\texttt{10.1109/IRASET68627.2026.11538879}. The entries that lacked volume and
pages were preprints typed as journal articles; they are now identified as
preprints with their arXiv identifiers, so no entry claims a volume it does not
have, and two IEEE entries are marked early access. We were unable to confirm
archival versions for the remaining preprints (Opacus, Flower, and the Green-AI
estimation tools); if the reviewer has a specific published version in mind for
[12], we will substitute it. Anguita et al.\ carried a spurious ``and others''
--- the paper has five authors --- and is now cited by its ESANN paper number,
since we could not confirm against the proceedings the page range that
circulates in secondary sources.

\textbf{6. Page limit.} The revision is six pages. Sections~V-A and V-B have
been condensed to the implications not already stated in Section~IV, and the
measurement-hazard discussion (Section~V-D) is preserved in full and extended
with the participation hazard of Point 3.

\end{document}
